\documentclass[aspectratio=169,9pt]{beamer}

% Theme and Colors
\usetheme{Madrid}
\usecolortheme{whale}

\definecolor{hcmutblue}{RGB}{0, 51, 102}
\definecolor{accentblue}{RGB}{37, 99, 235}
\definecolor{darkslate}{RGB}{30, 41, 59}
\definecolor{cardbg}{RGB}{241, 245, 249}

\setbeamercolor{palette primary}{bg=hcmutblue,fg=white}
\setbeamercolor{palette secondary}{bg=accentblue,fg=white}
\setbeamercolor{palette tertiary}{bg=hcmutblue!90,fg=white}
\setbeamercolor{titlelike}{bg=hcmutblue,fg=white}
\setbeamercolor{structure}{fg=hcmutblue}
\setbeamercolor{block title}{bg=hcmutblue,fg=white}
\setbeamercolor{block body}{bg=cardbg,fg=darkslate}
\setbeamercolor{block title example}{bg=accentblue,fg=white}
\setbeamercolor{block body example}{bg=cardbg,fg=darkslate}
\setbeamercolor{block title alert}{bg=red!80!black,fg=white}
\setbeamercolor{block body alert}{bg=red!5!white,fg=darkslate}

% Packages
\usepackage[utf8]{inputenc}
\usepackage[vietnamese]{babel}
\usepackage{booktabs}
\usepackage{tabularx}
\usepackage{graphicx}
\usepackage{amsmath,amssymb}
\usepackage{microtype}
\usepackage{pifont}

% Footer configuration
\setbeamertemplate{navigation symbols}{}
\setbeamertemplate{footline}{
  \leavevmode%
  \hbox{%
  \begin{beamercolorbox}[wd=.333333\paperwidth,ht=2.25ex,dp=1ex,center]{author in head/foot}%
    \usebeamerfont{author in head/foot}CO3021 Group 02 (HCMUT)
  \end{beamercolorbox}%
  \begin{beamercolorbox}[wd=.333333\paperwidth,ht=2.25ex,dp=1ex,center]{title in head/foot}%
    \usebeamerfont{title in head/foot}CO3021 $\cdot$ Group 02 | Other Indexing Structures
  \end{beamercolorbox}%
  \begin{beamercolorbox}[wd=.333333\paperwidth,ht=2.25ex,dp=1ex,right]{date in head/foot}%
    \usebeamerfont{date in head/foot}Tháng 10/2026\hspace*{2em}
    \insertframenumber{} / \inserttotalframenumber\hspace*{2ex} 
  \end{beamercolorbox}}%
  \vskip0pt%
}

% Title Metadata
\title[Other Indexing Structures]{\textbf{Other Indexing Structures and Indexes in Today's DBMSs}}
\subtitle{Nghiên cứu Thực nghiệm theo Workload trên Microsoft SQL Server 2025 \& MongoDB 7.0}
\author[Nhóm 02]{Nhóm 02 --- Giảng viên hướng dẫn: TS. Phan Trọng Nhân}
\institute[HCMUT]{Trường Đại học Bách khoa TP.HCM --- Khoa Khoa học \& Kỹ thuật Máy tính}
\date{Tháng 10/2026}

\begin{document}

% -----------------------------------------------------------------------------
% SLIDE 1: TITLE SLIDE
% -----------------------------------------------------------------------------
\begin{frame}
  \titlepage
\end{frame}

% -----------------------------------------------------------------------------
% SLIDE 2: GROUP ROSTER & RESPONSIBILITIES
% -----------------------------------------------------------------------------
\begin{frame}{Danh Sách Thành Viên \& Phân Công Nhiệm Vụ}
\begin{table}
\centering\footnotesize
\begin{tabularx}{\linewidth}{@{}l l X@{}}
\toprule
\textbf{MSSV} & \textbf{Họ và Tên} & \textbf{Trách nhiệm và Phân công chính trong đồ án} \\
\midrule
2252396 & Nguyễn Trung Kiên & Trưởng nhóm, kiến trúc tổng thể, benchmark UC1 Compound B-Tree \\
2353196 & Triệu Thuận Toàn & Thực nghiệm UC2 Filtered / Partial Index trên SQL/MongoDB \\
2352988 & Hồ Lê Minh Quân & Thực nghiệm UC3 Covering Index (INCLUDE vs Compound Keys) \\
2352608 & Huỳnh Minh Khởi & Thực nghiệm UC4 Full-Text / Text Search \& Inverted Index \\
2212828 & Vũ Minh Quân & Thực nghiệm UC5 Clustered Columnstore Index (SQL Server 1M) \\
2352951 & Lâm Thanh Phương & Thực nghiệm UC6 Multikey vs Normalized Child Index \\
2352699 & Thái Thiên Hải Long & Thiết kế dữ liệu EDA, kiểm định toàn vẹn dữ liệu Phase 1/Phase 8 \\
2352708 & Đinh Cao Thiên Lộc & Phân tích Actual Execution Plans và toán tử vật lý \\
2352967 & Ngô Phương Quang & Đo lường chi phí ghi, build time và freshness delay \\
2352758 & Phan Đức Minh & Xây dựng Web Demo Streamlit và báo cáo LaTeX V3 \\
\bottomrule
\end{tabularx}
\end{table}
\end{frame}

% -----------------------------------------------------------------------------
% SLIDE 3: AGENDA
% -----------------------------------------------------------------------------
\begin{frame}{Nội Dung Báo Cáo Thuyết Trình}
\begin{columns}[T]
\begin{column}{0.48\textwidth}
\begin{block}{Phần I: Đặt Vấn Đề \& Phương Pháp Luận}
\begin{itemize}
  \item \textbf{Chương 1:} 3 Câu hỏi cốt lõi \& Phân loại chỉ mục
  \item \textbf{Chương 2:} Thiết kế 6 Use Case theo hướng Workload-first
  \item \textbf{Chương 3:} 3 Bộ dữ liệu EDA \& 5 Quy chuẩn đo lường
\end{itemize}
\end{block}
\begin{block}{Phần II: Kết Quả Thực Nghiệm Sáu Use Case}
\begin{itemize}
  \item \textbf{UC1:} Compound B-Tree Index (Lọc đa cột \& Sắp xếp)
  \item \textbf{UC2:} Filtered / Partial Index (Dữ liệu lệch $0.24\%$)
  \item \textbf{UC3:} Covering Index (Triệt tiêu RID/Doc Lookup)
\end{itemize}
\end{block}
\end{column}

\begin{column}{0.48\textwidth}
\begin{block}{Phần II (tiếp): Các Cấu Trúc Khác B-Tree}
\begin{itemize}
  \item \textbf{UC4:} Full-Text / Inverted Index \& Bẫy ngữ nghĩa
  \item \textbf{UC5:} Clustered Columnstore (Nén 23:1, Batch Mode)
  \item \textbf{UC6:} Multikey vs. Normalized Child Index (Tags)
\end{itemize}
\end{block}
\begin{block}{Phần III: Tổng Hợp \& Trực Quan Hóa}
\begin{itemize}
  \item \textbf{Chương 5:} Bảng Quyết Định Kỹ Thuật Tổng Hợp
  \item \textbf{4 Cạm bẫy} phương pháp luận cần tránh
  \item \textbf{Live Demo:} Giới thiệu Web Streamlit tương tác
\end{itemize}
\end{block}
\end{column}
\end{columns}
\end{frame}

% -----------------------------------------------------------------------------
% SLIDE 4: CHAPTER 1 - PROBLEM & 3 CORE QUESTIONS
% -----------------------------------------------------------------------------
\begin{frame}{Vấn Đề Nghiên Cứu \& Ba Câu Hỏi Cốt Lõi}
\begin{block}{Mục Tiêu Đồ Án Môn Học}
Nghiên cứu các cấu trúc chỉ mục mở rộng \emph{ngoài slide lý thuyết} và các phép đánh chỉ mục hiện đại trong DBMS thực tế (Microsoft SQL Server 2025 \& MongoDB 7.0), trả lời dứt khoát 3 câu hỏi của giảng viên:
\end{block}

\begin{columns}[T]
\begin{column}{0.32\textwidth}
\begin{block}{\ding{182} What are they?}
\footnotesize
Các cấu trúc vật lý và kỹ thuật index này là gì? Bản chất lưu trữ trên đĩa, bộ nhớ đệm và cấu trúc dữ liệu bên dưới là gì?
\end{block}
\end{column}

\begin{column}{0.32\textwidth}
\begin{block}{\ding{183} How are they supported?}
\footnotesize
DBMS hiện thực và tối ưu hóa chúng như thế nào? (Plan cache, Seek predicate, Batch Mode, Background Catalog).
\end{block}
\end{column}

\begin{column}{0.32\textwidth}
\begin{alertblock}{\ding{184} How they support?}
\footnotesize
Chúng cải thiện hiệu năng cho \textbf{workload thực tế} nào? Đạt tốc độ bao nhiêu và phải trả giá bằng \textbf{đánh đổi gì}?
\end{alertblock}
\end{column}
\end{columns}

\vspace{0.4cm}
\centering
\textbf{\color{red!80!black}\ding{52} Tránh bẫy nghiên cứu: Không tạo index ``cho vui'', mọi index đều phải phục vụ một Use Case rõ ràng!}
\end{frame}

% -----------------------------------------------------------------------------
% SLIDE 5: TAXONOMY - PHYSICAL STRUCTURES VS B-TREE TECHNIQUES
% -----------------------------------------------------------------------------
\begin{frame}{Phân Biệt Cấu Trúc Vật Lý vs. Kỹ Thuật Thiết Kế B-Tree}
\begin{columns}[T]
\begin{column}{0.48\textwidth}
\begin{block}{Cấu Trúc Vật Lý Cơ Sở (Physical Storage)}
Định nghĩa \textbf{cách thức tổ chức các byte trên đĩa/RAM}:
\begin{itemize}
  \item \textbf{Heap / Clustered Rowstore:} Tổ chức theo từng trang (Page 8KB trong SQL, WiredTiger pages trong MongoDB).
  \item \textbf{Inverted Index (UC4):} Tách từ (Word Breaker), tạo bảng từ vựng (Lexicon) và Posting Lists trỏ tới tài liệu.
  \item \textbf{Clustered Columnstore (UC5):} Chia bảng thành rowgroup, nén dữ liệu theo từng cột (Column Segment) với từ điển và RLE.
\end{itemize}
\end{block}
\end{column}

\begin{column}{0.48\textwidth}
\begin{exampleblock}{Kỹ Thuật Thiết Kế Trên Cây B-Tree}
Khai thác \textbf{cấu trúc B-Tree sẵn có} để tối ưu đường truy cập:
\begin{itemize}
  \item \textbf{Compound B-Tree (UC1):} Ghép nhiều cột theo thứ tự từ điển; tối ưu lọc đa chiều và triệt tiêu toán tử Sort.
  \item \textbf{Filtered / Partial Index (UC2):} B-Tree chỉ lưu trữ tập con bản ghi thỏa mãn biểu thức vị từ (Predicate).
  \item \textbf{Covering Index (UC3):} Bổ sung cột phụ trợ vào trang lá (\texttt{INCLUDE} hoặc Compound) để triệt tiêu Base Lookups.
  \item \textbf{Multikey Index (UC6):} Nhân bản các khóa trỏ cho từng phần tử trong mảng (NoSQL).
\end{itemize}
\end{exampleblock}
\end{column}
\end{columns}
\end{frame}

% -----------------------------------------------------------------------------
% SLIDE 6: DATASETS & EDA
% -----------------------------------------------------------------------------
\begin{frame}{Ba Bộ Dữ Liệu Thực Nghiệm \& Đặc Trưng EDA}
\begin{columns}[T]
\begin{column}{0.48\textwidth}
\begin{block}{Ba Bộ Dữ Liệu Thực Nghiệm}
\footnotesize
\begin{itemize}
  \item \textbf{BIG\_TABLE (100K \& 1M dòng):} 14 cột, kế thừa từ Oracle \texttt{ALL\_OBJECTS}. Dùng cho UC1, UC2, UC3, UC5.
  \item \textbf{Corpus Tin Tức (10,000 bài báo):} Văn bản tiếng Anh thực tế (CNN/Reuters), cột \texttt{BODY} dài trung vị 371 từ (max 900+ từ). Dùng cho UC4 Full-Text.
  \item \textbf{Corpus Mảng Tags (10,000 thực thể):} 30,471 quan hệ tag (mỗi doc 3--5 tags). Dùng cho UC6 Multikey.
\end{itemize}
\end{block}
\end{column}

\begin{column}{0.48\textwidth}
\begin{exampleblock}{Đặc Tính Phân Phối Dữ Liệu (EDA)}
\footnotesize
\begin{itemize}
  \item \textbf{Độ lệch cực đoan (UC2):} Cột \texttt{TEMPORARY = 'Y'} chỉ chiếm đúng \textbf{$0.24\%$} (240 dòng ở 100K; 2,405 dòng ở 1M), $99.76\%$ là 'N'.
  \item \textbf{Độ lệch cột lọc (UC1):} Cột \texttt{OWNER = 'SYS'} chiếm tới \textbf{$71.8\%$}, trong khi các giá trị khác rất thưa.
  \item \textbf{Tính toàn vẹn (Integrity):} SHA-256 đối chiếu giữa SQL Server và MongoDB trùng khớp 100\% tập dữ liệu thô.
\end{itemize}
\end{exampleblock}
\end{column}
\end{columns}
\end{frame}

% -----------------------------------------------------------------------------
% SLIDE 7: FAIRNESS CONSTRAINTS & WORKFLOW
% -----------------------------------------------------------------------------
\begin{frame}{Năm Quy Chuẩn Thực Nghiệm Công Bằng (Fairness Rules)}
\begin{block}{Quy Trình Đo Đạc \& Giám Sát Khoa Học}
\begin{enumerate}
  \item \textbf{Disposable Working Copies (Bản sao dùng một lần):} Mỗi thí nghiệm clone một bản sao sạch (\texttt{UC*\_Read} và \texttt{UC*\_Write}) và tước bỏ index cũ, đảm bảo các kịch bản ghi không làm sai lệch số dòng đọc.
  \item \textbf{Warmups \& 25 Timed Runs:} Chạy 2 lần warmup không tính giờ, sau đó đo 25 lần timed runs, vật chất hóa toàn bộ dòng và lấy trung vị (p50) làm chỉ số chính; ghi nhận p95 để kiểm soát đuôi độ trễ.
  \item \textbf{Separate Actual Execution Plan Capture (Bắt kế hoạch riêng biệt):} Không bật cờ trace trong lúc đo giờ; sau 25 lần đo mới chạy lượt chẩn đoán riêng (\texttt{SET STATISTICS IO, XML ON} trong SQL; \texttt{explain('executionStats')} trong MongoDB).
  \item \textbf{Đo lường đa chiều: Storage, Build Time \& Write Maintenance:} Đo kích thước index trên đĩa (\texttt{sys.dm\_db\_partition\_stats} / \texttt{indexSizes}), thời gian build (ms), thông lượng batch INSERT (rows/s) và độ trễ UPDATE.
  \item \textbf{Strict Intra-DBMS Comparison Rule (Chỉ so sánh nội bộ cùng DBMS):} SQL Server dùng server-side batch timer; MongoDB đo client cursor. Tuyệt đối chỉ so sánh A vs B vs C trong cùng DBMS, không so ms chéo giữa 2 hệ thống!
\end{enumerate}
\end{block}
\end{frame}

% -----------------------------------------------------------------------------
% SLIDE 8: 6-UC OVERVIEW TABLE
% -----------------------------------------------------------------------------
\begin{frame}{Tổng Quan Sáu Use Case Thực Nghiệm (UC1--UC6)}
\begin{table}
\centering\footnotesize
\begin{tabularx}{\linewidth}{@{}l l l X@{}}
\toprule
\textbf{UC} & \textbf{Bài Toán Workload} & \textbf{Chỉ Mục Giải Pháp} & \textbf{Mở Rộng Lý Thuyết} \\
\midrule
\textbf{UC1} & Lọc nhiều cột \& sắp xếp & \textbf{Compound B-Tree} & Thứ tự từ điển, quy tắc tiền tố trái, triệt tiêu toán tử Sort. \\
\textbf{UC2} & Cờ trạng thái hiếm ($0.24\%$ Y) & \textbf{Filtered / Partial Index} & Chỉ mục có vị từ lọc, tiết kiệm $>98\%$ đĩa, giới hạn với ngoài tập. \\
\textbf{UC3} & Range query lấy thêm 4 cột & \textbf{Covering Index} (INCLUDE) & Lưu payload ở trang lá, loại bỏ $100\%$ RID Lookup / FETCH. \\
\textbf{UC4} & Tìm từ/cụm từ trong bài báo & \textbf{Full-Text / Text Index} & Chỉ mục đảo, bẫy ngữ nghĩa (khác ID), Async lag vs Sync write penalty. \\
\textbf{UC5} & Analytical GROUP BY 1M dòng & \textbf{Clustered Columnstore} (CCI) & Lưu trữ cột nén (23:1), Batch Mode execution, Delta store writes. \\
\textbf{UC6} & Đối tượng chứa 1 hoặc 2 tags & \textbf{Multikey vs. Child Index} & Nhân bản index entry mảng trong NoSQL vs. bảng con chuẩn hóa 1NF. \\
\bottomrule
\end{tabularx}
\end{table}
\end{frame}

% -----------------------------------------------------------------------------
% SLIDE 9: UC1 COMPOUND B-TREE (SETUP & 4 QUERIES)
% -----------------------------------------------------------------------------
\begin{frame}{UC1: Multi-Column Lookups \& Sorting $\to$ Compound B-Tree}
\textbf{Bài toán thực tế:} Lọc kết hợp \texttt{OWNER}, \texttt{OBJECT\_TYPE} và sắp xếp theo \texttt{OBJECT\_ID}. Single-column index trên \texttt{OWNER} chỉ thu hẹp được một phần, sau đó DBMS phải quét và lọc phần còn lại.

\begin{block}{Thiết Kế Bốn Câu Truy Vấn Độc Lập Để Cách Ly Từng Cơ Chế}
\footnotesize
\begin{itemize}
  \item \textbf{Query U1a (2 dòng khớp):} \texttt{WHERE OWNER='SYS' AND OBJECT\_TYPE='TABLE' AND OBJECT\_ID BETWEEN 100 AND 200} $\to$ Kiểm tra trường hợp lý tưởng nhất khi cả 3 điều kiện khớp hoàn hảo với tiền tố index.
  \item \textbf{Query U1b (3,405 dòng khớp):} \texttt{WHERE OWNER='SYS' AND OBJECT\_TYPE='TABLE'} (Không có dải ID) $\to$ Đánh giá chi phí RID Lookup khi số dòng kết quả lớn (Tipping Point).
  \item \textbf{Query U1c (2 dòng khớp):} \texttt{WHERE OBJECT\_TYPE='TABLE' AND OBJECT\_ID BETWEEN 100 AND 200} $\to$ \textbf{Vi phạm tiền tố trái (Skip Leading Key)}: Kiểm tra optimizer có dùng được index không.
  \item \textbf{Query U1d (398 dòng khớp):} \texttt{WHERE OWNER='SYS' AND OBJECT\_TYPE='VIEW' ORDER BY OBJECT\_ID} $\to$ Kiểm tra khả năng \textbf{triệt tiêu toán tử Sort (Sort Elimination)}.
\end{itemize}
\end{block}
\end{frame}

% -----------------------------------------------------------------------------
% SLIDE 10: UC1 RESULTS & EXECUTION PLAN
% -----------------------------------------------------------------------------
\begin{frame}{UC1: Kết Quả Đo Đạc, Execution Plan \& Đánh Đổi}
\begin{table}
\centering\scriptsize
\begin{tabularx}{\linewidth}{@{}l l r r r r X@{}}
\toprule
\textbf{DBMS} & \textbf{Truy Vấn} & \textbf{A p50} & \textbf{B p50} & \textbf{C p50} & \textbf{Speedup} & \textbf{Hành Vi Actual Execution Plan Nổi Bật} \\
\midrule
SQL Server & \texttt{U1a} (2 rows) & 32.36 ms & 24.53 ms & \textbf{9.46 ms} & \textbf{3.42$\times$} & Index Seek trỏ đúng 2 dòng, giảm reads từ 2,730 $\to$ 2 trang \\
SQL Server & \texttt{U1b} (3,405 rows) & 104.48 ms & 84.29 ms & 97.99 ms & 1.07$\times$ & Bỏ qua C để Table Scan vì 3,405 RID lookups quá đắt \\
SQL Server & \texttt{U1c} (2 rows) & 37.01 ms & 31.97 ms & \textbf{20.04 ms} & \textbf{1.85$\times$} & Thiếu leading key, phải quét toàn bộ cây index thay vì Seek \\
SQL Server & \texttt{U1d} (398 rows) & 38.02 ms & 11.12 ms & \textbf{10.52 ms} & \textbf{3.61$\times$} & \textbf{Triệt tiêu Sort} hoàn toàn nhờ cây index đã sắp sẵn \\
\midrule
MongoDB & \texttt{U1a} (2 docs) & 241.03 ms & 462.78 ms & \textbf{2.34 ms} & \textbf{103.0$\times$} & B quét 71K keys + 71K fetch docs; C seek đúng 2 keys \\
MongoDB & \texttt{U1b} (3,405 docs) & 253.16 ms & 457.73 ms & \textbf{39.31 ms} & \textbf{6.44$\times$} & B sụp đổ hiệu năng (chậm hơn A 2x); C gom cụm seek hẹp \\
MongoDB & \texttt{U1c} (2 docs) & 228.33 ms & 204.50 ms & 232.18 ms & 0.98$\times$ & Không có leading key $\to$ Mongo buộc phải COLLSCAN \\
MongoDB & \texttt{U1d} (398 docs) & 236.43 ms & 10.04 ms & \textbf{7.97 ms} & \textbf{29.7$\times$} & Trả kết quả ngay không cần toán tử bộ nhớ SORT \\
\bottomrule
\end{tabularx}
\end{table}

\vspace{-0.2cm}
\begin{columns}[T]
\begin{column}{0.48\textwidth}
\begin{block}{Giải Thích Hiện Tượng Actual Plan}
\scriptsize
\begin{itemize}
  \item \textbf{Mongo B chậm hơn A gần 2x trên U1a/U1b:} Do \texttt{OWNER='SYS'} duyệt 71,802 key và FETCH 71,802 docs ngẫu nhiên từ RAM/đĩa! Compound C thu hẹp seek về đúng 2 docs.
  \item \textbf{SQL Server bỏ C trên U1b:} 3,405 lần RID Lookup ngẫu nhiên đắt hơn quét tuần tự 2,730 trang heap $\to$ Optimizer tự động chọn \texttt{Table Scan}.
\end{itemize}
\end{block}
\end{column}

\begin{column}{0.48\textwidth}
\begin{alertblock}{Đánh Đổi Dung Lượng \& Chi Phí Ghi (Table B)}
\scriptsize
\begin{itemize}
  \item \textbf{Dung lượng index:} Compound C lớn hơn B từ $+88\%$ đến $+96\%$ (29.2 MB vs 15.5 MB ở SQL 1M).
  \item \textbf{Thông lượng chèn:} INSERT trên Mongo giảm $-33.5\%$ (11,628 $\to$ 7,736 docs/s).
  \item \textbf{Kết luận:} Chỉ nên tạo khi lọc có leading key hoặc ORDER BY khớp cây.
\end{itemize}
\end{alertblock}
\end{column}
\end{columns}
\end{frame}

% -----------------------------------------------------------------------------
% SLIDE 11: UC2 FILTERED & PARTIAL INDEX (SETUP)
% -----------------------------------------------------------------------------
\begin{frame}{UC2: Filtered \& Partial Index (Dữ Liệu Lệch Cực Đoan)}
\textbf{Bài toán thực tế:} Cờ \texttt{TEMPORARY = 'Y'} chỉ chiếm $0.24\%$ (2,405 dòng ở 1M), còn $99.76\%$ là 'N'. Dashboard liên tục truy vấn dòng 'Y', trong khi các tác vụ batch hiếm khi quét toàn bộ. Full index lãng phí $>99\%$ bộ nhớ vào các dòng 'N'.

\begin{block}{Thiết Kế Ba Truy Vấn Kiểm Thử Hành Vi Bộ Tối Ưu Hóa}
\footnotesize
\begin{enumerate}
  \item \textbf{Query Y (Inside-Subset --- 615 dòng ở 1M):} \texttt{WHERE TEMPORARY = 'Y' AND OBJECT\_TYPE = 'TABLE' AND OBJECT\_ID BETWEEN 1 AND 10000} $\to$ Kiểm tra hiệu năng đọc khi truy vấn đúng vào tập con được index lọc bảo vệ.
  \item \textbf{Query N (Outside-Subset --- 15,810 dòng ở 1M):} \texttt{WHERE TEMPORARY = 'N' AND ...} $\to$ Kiểm tra phản ứng của optimizer khi dữ liệu không có trong index C.
  \item \textbf{Query @p (Parameterized Probe trong SQL Server):} \texttt{WHERE TEMPORARY = @p} $\to$ Chứng minh cơ chế plan caching khiến optimizer buộc phải từ chối Filtered Index nếu không biên dịch lại.
\end{enumerate}
\end{block}

\begin{exampleblock}{Cơ Chế Chỉ Mục Lọc / Bán Phần}
\footnotesize
SQL Server: \texttt{CREATE INDEX ... WHERE TEMPORARY = 'Y'}; MongoDB: \texttt{partialFilterExpression: \{ temporary: "Y" \}}. Cây B-Tree chỉ lưu trữ đúng 2,405 dòng thay vì 1,000,000 dòng!
\end{exampleblock}
\end{frame}

% -----------------------------------------------------------------------------
% SLIDE 12: UC2 RESULTS & STORAGE SAVINGS
% -----------------------------------------------------------------------------
\begin{frame}{UC2: Kết Quả Đo Đạc \& Mức Tiết Kiệm Dung Lượng $>99.5\%$}
\begin{table}
\centering\footnotesize
\begin{tabularx}{\linewidth}{@{}l l r r r X@{}}
\toprule
\textbf{Hệ Quản Trị} & \textbf{Tập Truy Vấn} & \textbf{Baseline A} & \textbf{Full B} & \textbf{Filtered/Partial C} & \textbf{Công Việc Vật Lý (A $\to$ B $\to$ C)} \\
\midrule
SQL Server 1M & Inside 'Y' (615 rows) & 92.32 ms & 9.00 ms & \textbf{5.99 ms} & $28,032 \to 81 \to \mathbf{6\text{ trang}}$ (\textbf{15.4$\times$}) \\
SQL Server 1M & Outside 'N' (15,810) & 347.95 ms & \textbf{292.28 ms} & 349.62 ms & $28,032 \to \mathbf{81} \to 28,032$ (Quét toàn bộ Heap) \\
\midrule
MongoDB 1M & Inside 'Y' (615 docs) & 2,418.9 ms & 17.17 ms & \textbf{6.13 ms} & 1M docs $\to 1,670 \to \mathbf{616\text{ keys}}$ (\textbf{395$\times$}) \\
MongoDB 1M & Outside 'N' (15,810) & 2,447.9 ms & \textbf{69.45 ms} & 2,438.7 ms & 1M docs $\to \mathbf{15,852} \to 1\text{M docs}$ (COLLSCAN) \\
\bottomrule
\end{tabularx}
\end{table}

\vspace{-0.2cm}
\begin{columns}[T]
\begin{column}{0.48\textwidth}
\begin{exampleblock}{Tiết Kiệm Dung Lượng Phi Thường (Table B)}
\scriptsize
\begin{itemize}
  \item \textbf{SQL Server (1M):} Full Index B: $43.74$\,MB $\to$ Filtered Index C chỉ chiếm \textbf{$0.13$\,MB (Tiết kiệm $-99.71\%$!)}.
  \item \textbf{MongoDB (1M):} Full Index B: $11.95$\,MB $\to$ Partial Index C chỉ chiếm \textbf{$0.05$\,MB (Tiết kiệm $-99.57\%$!)}.
  \item \textbf{Thời gian build:} Dựng C nhanh gấp 3.4x--9.7x so với B (179 ms vs 1,741 ms ở SQL).
\end{itemize}
\end{exampleblock}
\end{column}

\begin{column}{0.48\textwidth}
\begin{alertblock}{Giới Hạn Bắt Buộc Cần Ghi Nhớ}
\scriptsize
\begin{itemize}
  \item \textbf{Ngoài tập lọc (Query N):} Cả 2 DBMS sụp đổ về Full Table Scan (349 ms / 2,438 ms).
  \item \textbf{Tham số hóa (@p trong SQL):} Optimizer từ chối Filtered Index nếu không dùng \texttt{RECOMPILE} để bảo đảm tính đúng đắn của Plan Cache.
\end{itemize}
\end{alertblock}
\end{column}
\end{columns}
\end{frame}

% -----------------------------------------------------------------------------
% SLIDE 13: UC3 COVERING INDEX (SETUP & TIERS)
% -----------------------------------------------------------------------------
\begin{frame}{UC3: Covering Index (INCLUDE vs. Compound Keys)}
\textbf{Bài toán thực tế:} Lọc theo dải \texttt{OBJECT\_ID} và trả về 4 cột mô tả (\texttt{ID, OWNER, OBJECT\_NAME, OBJECT\_TYPE}). Chỉ mục hẹp (Narrow B) tìm được con trỏ nhưng phải thực hiện hàng ngàn lần nhảy ngẫu nhiên (\texttt{RID Lookup} / \texttt{FETCH}) về bảng gốc.

\begin{block}{Thiết Kế Ba Tầng Độ Chọn Lọc (Selectivity Tiers)}
\footnotesize
\begin{enumerate}
  \item \textbf{Point Lookup (1 dòng):} \texttt{WHERE OBJECT\_ID = 16 AND ID = 1} $\to$ Cả Narrow B và Covering C đều hoàn thành trong $\le 3.5$ ms. Phủ index không đem lại lợi ích rõ rệt.
  \item \textbf{Small Range (60--98 dòng):} \texttt{WHERE OBJECT\_ID BETWEEN 1 AND 50} (100K) hoặc \texttt{1 AND 5} (1M) $\to$ Dữ liệu nằm sẵn trong buffer pool, chi phí lookup của B vẫn rất nhỏ.
  \item \textbf{Medium Range (1,694--2,775 dòng):} \texttt{BETWEEN 1 AND 1000} (100K) hoặc \texttt{1 AND 200} (1M) $\to$ \textbf{Điểm bùng nổ (Tipping Point)}: Số lần RID Lookup quá lớn khiến Narrow B bị sụp đổ hiệu năng, Covering C phát huy sức mạnh tối đa!
\end{enumerate}
\end{block}

\begin{exampleblock}{Cơ Chế Phủ Chỉ Mục}
\footnotesize
SQL Server dùng mệnh đề \texttt{INCLUDE (OWNER, OBJECT\_NAME, OBJECT\_TYPE)} lưu payload tại trang lá; MongoDB đưa toàn bộ 4 trường vào khóa Compound và gọi phép chiếu \texttt{\{ \_id: 0 \}} để đạt \texttt{PROJECTION\_COVERED}.
\end{exampleblock}
\end{frame}

% -----------------------------------------------------------------------------
% SLIDE 14: UC3 RESULTS & BASE LOOKUP ELIMINATION
% -----------------------------------------------------------------------------
\begin{frame}{UC3: Kết Quả Thực Nghiệm, Triệt Tiêu Lookup \& Đánh Đổi}
\begin{table}
\centering\footnotesize
\begin{tabularx}{\linewidth}{@{}l l r r r X@{}}
\toprule
\textbf{Hệ Quản Trị} & \textbf{Tầng Chọn Lọc} & \textbf{Baseline A} & \textbf{Narrow B} & \textbf{Covering C} & \textbf{Công Việc Vật Lý (B $\to$ C)} \\
\midrule
SQL Server 100K & Point (1 row) & 26.09 ms & 0.00$^*$ ms & 0.00$^*$ ms & $4 \to \mathbf{2\text{ trang}}$ (Triệt tiêu RID Lookup) \\
SQL Server 100K & Medium (1,694 rows) & 49.59 ms & 52.22 ms & \textbf{22.63 ms} & 2,730 (Scan) $\to \mathbf{29\text{ trang}}$ (\textbf{2.31$\times$}) \\
SQL Server 1M & Medium (2,775 rows) & 113.93 ms & 49.78 ms & \textbf{32.23 ms} & $2,785 \to \mathbf{49\text{ trang}}$ (\textbf{1.54$\times$}) \\
\midrule
MongoDB 100K & Small (98 docs) & 225.39 ms & \textbf{2.83 ms} & 4.40 ms & 98 docs fetched $\to \mathbf{0\text{ docs}}$ (C chậm hơn do scan payload) \\
MongoDB 100K & Medium (1,694 docs) & 240.47 ms & 17.21 ms & \textbf{14.64 ms} & 1,694 docs $\to \mathbf{0\text{ docs}}$ (\textbf{PROJECTION\_COVERED}) \\
MongoDB 1M & Medium (2,775 docs) & 2,579.4 ms & 23.91 ms & \textbf{20.12 ms} & 2,775 docs $\to \mathbf{0\text{ docs}}$ (\textbf{128$\times$} vs Baseline A) \\
\bottomrule
\end{tabularx}
\end{table}

\vspace{-0.2cm}
\begin{columns}[T]
\begin{column}{0.48\textwidth}
\begin{block}{Cơ Chế \& Tipping Point Bùng Nổ}
\scriptsize
\begin{itemize}
  \item \textbf{Triệt tiêu 100\% Base Lookups:} SQL C không còn toán tử \texttt{RID Lookup}; Mongo C đạt stage \texttt{PROJECTION\_COVERED} với \texttt{totalDocsExamined = 0}.
  \item \textbf{Tipping Point ở 100K Medium:} Chi phí 1,694 lần lookup khiến SQL B bỏ index chọn Table Scan; C chỉ đọc 29 trang lá (nhanh hơn 2.31x).
\end{itemize}
\end{block}
\end{column}

\begin{column}{0.48\textwidth}
\begin{alertblock}{Đánh Đổi Phình To Index (Table B)}
\scriptsize
\begin{itemize}
  \item \textbf{Phình to dung lượng (1M):} Index C tăng \textbf{5.94x trong SQL} (127.4 MB) và \textbf{10.47x trong Mongo} (66.9 MB).
  \item \textbf{Chi phí ghi:} INSERT throughput giảm $-14.7\%$ trên SQL Server.
  \item \textbf{Chỉ định dùng:} Chỉ nên áp dụng khi range query có tần suất thực thi cao!
\end{itemize}
\end{alertblock}
\end{column}
\end{columns}
\end{frame}

% -----------------------------------------------------------------------------
% SLIDE 15: UC4 FULL-TEXT / TEXT INDEX (SETUP & 6 QUERIES)
% -----------------------------------------------------------------------------
\begin{frame}{UC4: Full-Text \& Text Search (Tìm Kiếm Trên Văn Bản Báo)}
\textbf{Bài toán thực tế:} Tìm kiếm từ/cụm từ bên trong 10,000 bài báo tiếng Anh (\texttt{BODY} trung vị 371 từ, dài tới 900+ từ). Cây B-Tree chuỗi thông thường vô dụng; quét chuỗi \texttt{LIKE '\%term\%'} hoặc \texttt{\$regex} mất từ 3 đến 4 giây CPU!

\begin{block}{Thiết Kế Sáu Câu Truy Vấn Cô Lập Độ Chọn Lọc}
\footnotesize
\begin{columns}[T]
\begin{column}{0.48\textwidth}
\textbf{Ba Truy Vấn Từ Đơn (Single Keywords):}
\begin{itemize}
  \item \texttt{election} (2,577 bài --- Tần suất cao)
  \item \texttt{climate} (326 bài --- Tần suất trung bình)
  \item \texttt{bitcoin} (7 bài --- Tần suất hiếm)
\end{itemize}
\end{column}
\begin{column}{0.48\textwidth}
\textbf{Ba Truy Vấn Cụm Từ Chính Xác (Exact Phrases):}
\begin{itemize}
  \item \texttt{"white house"} (1,851 bài --- Phổ biến)
  \item \texttt{"climate change"} (214 bài --- Trung bình)
  \item \texttt{"stock market"} (45 bài --- Ít)
\end{itemize}
\end{column}
\end{columns}
\end{block}

\begin{exampleblock}{Cơ Chế Chỉ Mục Đảo (Inverted Index)}
\footnotesize
Tách từ (Word Breaker) $\to$ Lọc Stop-words (the, is, and) $\to$ Gốc từ (Stemmer: elections $\to$ elect) $\to$ \textbf{Posting Lists} trỏ tới danh sách DOC\_ID.
\end{exampleblock}
\end{frame}

% -----------------------------------------------------------------------------
% SLIDE 16: UC4 RESULTS & SEMANTIC PARITY TRAP
% -----------------------------------------------------------------------------
\begin{frame}{UC4: Kết Quả Đo Đạc \& Bẫy Ngữ Nghĩa (Semantic Parity)}
\begin{columns}[T]
\begin{column}{0.48\textwidth}
\textbf{Microsoft SQL Server (LIKE vs. CONTAINS)}
\begin{table}
\centering\scriptsize
\resizebox{\linewidth}{!}{
\begin{tabular}{@{}l r r c r r@{}}
\toprule
Từ Khóa & A Dòng & C Dòng & Trùng ID? & A p50 & C p50 \\
\midrule
\texttt{election} & 2,577 & 2,311 & \texttt{No} & 3,294 ms & \textbf{21.0 ms} \\
\texttt{climate} & 326 & 324 & \texttt{No} & 3,983 ms & \textbf{17.4 ms} \\
\texttt{bitcoin} & 7 & 6 & \texttt{No} & 3,987 ms & \textbf{14.9 ms} \\
\texttt{"white house"} & 1,851 & 1,845 & \texttt{No} & 3,440 ms & \textbf{22.5 ms} \\
\texttt{"climate chg"} & 214 & 215 & \texttt{No} & 4,042 ms & \textbf{16.6 ms} \\
\texttt{"stock mkt"} & 45 & 37 & \texttt{No} & 4,060 ms & \textbf{14.1 ms} \\
\bottomrule
\end{tabular}
}
\end{table}
\begin{alertblock}{Bẫy Ngữ Nghĩa Trên SQL Server}
\scriptsize
Cả 6 query đều trả về \textbf{khác tập DOC\_ID}! LIKE tìm chuỗi con ('re-election'), CONTAINS tìm từ ngữ nghĩa. Đây là 2 câu hỏi nghiệp vụ khác nhau!
\end{alertblock}
\end{column}

\begin{column}{0.48\textwidth}
\textbf{MongoDB 7.0 (\$regex vs. \$text)}
\begin{table}
\centering\scriptsize
\resizebox{\linewidth}{!}{
\begin{tabular}{@{}l r r c r r@{}}
\toprule
Từ Khóa & A Doc & C Doc & Trùng ID? & A p50 & C p50 \\
\midrule
\texttt{election} & 2,577 & 3,083 & \texttt{No} & 384.6 ms & \textbf{15.2 ms} \\
\texttt{climate} & 326 & 326 & \textbf{\texttt{Yes}} & 299.9 ms & \textbf{6.4 ms} (\textbf{46.8$\times$}) \\
\texttt{bitcoin} & 7 & 7 & \textbf{\texttt{Yes}} & 197.5 ms & \textbf{1.8 ms} (\textbf{107.8$\times$}) \\
\texttt{"white house"} & 1,851 & 1,851 & \textbf{\texttt{Yes}} & 140.5 ms & \textbf{42.6 ms} (\textbf{3.3$\times$}) \\
\texttt{"climate chg"} & 214 & 214 & \textbf{\texttt{Yes}} & 288.4 ms & \textbf{52.4 ms} (\textbf{5.5$\times$}) \\
\texttt{"stock mkt"} & 45 & 45 & \textbf{\texttt{Yes}} & 419.6 ms & \textbf{25.6 ms} (\textbf{16.4$\times$}) \\
\bottomrule
\end{tabular}
}
\end{table}
\begin{exampleblock}{Kết Quả Mongo \& Stemming}
\scriptsize
5 query trả đúng $100\%$ ID có tốc độ tăng từ \textbf{3.3x đến 107.8x}! 'election' trên C trả 3,083 docs do quy về gốc từ 'elect', 'elected'.
\end{exampleblock}
\end{column}
\end{columns}
\end{frame}

% -----------------------------------------------------------------------------
% SLIDE 17: UC4 ARCHITECTURAL TRADEOFF (ASYNC VS SYNC)
% -----------------------------------------------------------------------------
\begin{frame}{UC4: Đánh Đổi Kiến Trúc: Bất Đồng Bộ (SQL) vs. Đồng Bộ (Mongo)}
\begin{columns}[T]
\begin{column}{0.48\textwidth}
\begin{block}{Microsoft SQL Server (Asynchronous Catalog)}
\footnotesize
\begin{itemize}
  \item \textbf{Cơ chế hoạt động:} Tiến trình ngầm (\texttt{fdhost.exe}) thu thập thay đổi qua crawl queue ngoài transaction engine.
  \item \textbf{Ưu điểm ghi:} INSERT throughput không bị ảnh hưởng trực tiếp lúc chèn dữ liệu ($14,500$ dòng/s).
  \item \textbf{Nhược điểm chí mạng (Freshness Lag):} Dòng mới chèn mất trung bình \textbf{$5.8$ giây} mới xuất hiện trong kết quả truy vấn \texttt{CONTAINS}!
  \item \textbf{Phù hợp:} Ứng dụng chấp nhận tính nhất quán cuối cùng (Eventual Consistency).
\end{itemize}
\end{block}
\end{column}

\begin{column}{0.48\textwidth}
\begin{alertblock}{MongoDB (Synchronous Inverted Index)}
\footnotesize
\begin{itemize}
  \item \textbf{Cơ chế hoạt động:} Mọi thao tác ghi văn bản đều phải tách từ (tokenize), băm gốc từ và cập nhật posting lists ngay trong transaction.
  \item \textbf{Ưu điểm đọc:} Dữ liệu có thể tìm kiếm được \textbf{ngay lập tức (Zero lag)} sau lệnh chèn.
  \item \textbf{Cái giá cực đắt cho thông lượng ghi:} INSERT throughput sụt giảm tới \textbf{$-95.3\%$} (từ 10,753 doc/s rớt xuống còn 509 doc/s)!
  \item \textbf{Phù hợp:} Ứng dụng đọc nhiều ghi ít, bắt buộc thấy dữ liệu tức thì.
\end{itemize}
\end{alertblock}
\end{column}
\end{columns}
\end{frame}

% -----------------------------------------------------------------------------
% SLIDE 18: UC5 CLUSTERED COLUMNSTORE INDEX (SETUP)
% -----------------------------------------------------------------------------
\begin{frame}{UC5: Clustered Columnstore Index (SQL Server 1M Dòng)}
\textbf{Bài toán thực tế:} Truy vấn phân tích \texttt{GROUP BY} trên 2 cột (\texttt{OBJECT\_TYPE, OBJECT\_ID}) trên bảng 14 cột. Rowstore Heap buộc phải nạp toàn bộ 219 MB (28,032 trang) vào bộ đệm để đọc cả 12 cột thừa.

\begin{block}{Cơ Chế Lưu Trữ Dạng Cột Nén (CCI) \& Hai Câu Truy Vấn}
\begin{itemize}
  \item Bảng được chia thành rowgroup ($1,048,576$ dòng). Mỗi cột được lưu thành \textbf{Column Segment} riêng biệt, nén sâu bằng từ điển và run-length encoding.
  \item Xử lý vector hóa theo mảng khoảng 900 dòng (\textbf{Batch Mode Execution}).
  \item \textbf{Truy vấn AGG (Analytical GROUP BY --- 35 nhóm):} \texttt{SELECT OBJECT\_TYPE, COUNT\_BIG(*), SUM(COALESCE(OBJECT\_ID,0)) FROM Table GROUP BY OBJECT\_TYPE}.
  \item \textbf{Truy vấn POINT (Unindexed Scan --- 1 dòng):} \texttt{SELECT ID, OWNER, OBJECT\_NAME, OBJECT\_TYPE FROM Table WHERE ID = 1} (Chạy khi cả hai bảng \textbf{không có B-tree trên ID}).
\end{itemize}
\end{block}
\end{frame}

% -----------------------------------------------------------------------------
% SLIDE 19: UC5 RESULTS & DELTA ROWGROUP
% -----------------------------------------------------------------------------
\begin{frame}{UC5: Kết Quả Đo Đạc, Tỉ Lệ Nén 23:1 \& OPEN Delta Rowgroup}
\begin{table}
\centering\footnotesize
\begin{tabularx}{\linewidth}{@{}l r r r X@{}}
\toprule
\textbf{Truy Vấn Nghiên Cứu} & \textbf{Rowstore Heap B} & \textbf{Clustered Columnstore C} & \textbf{Tỉ Lệ Tăng Tốc} & \textbf{Công Việc Vật Lý (Execution Plan)} \\
\midrule
\texttt{AGG} (GROUP BY 1M dòng) & 123.2 ms & \textbf{42.4 ms} & \textbf{2.90$\times$} & 28,032 trang Row Mode $\to$ \textbf{1,105 LOB reads Batch Mode} \\
\texttt{POINT} (Scan 1 dòng) & 88.8 ms & \textbf{3.6 ms} & \textbf{24.49$\times$} & 28,032 trang $\to$ 776 LOB reads (\emph{vẫn là scan, không phải seek}) \\
\bottomrule
\end{tabularx}
\end{table}

\vspace{-0.2cm}
\begin{columns}[T]
\begin{column}{0.48\textwidth}
\begin{exampleblock}{Tỉ Lệ Nén Vật Lý Phi Thường (Table B)}
\scriptsize
\begin{itemize}
  \item Rowstore Heap B: $219.00$\,MB ($28,032$ trang dữ liệu)
  \item Clustered Columnstore C: \textbf{$9.46$\,MB ($1,211$ trang)}
  \item \textbf{Đạt tỉ lệ nén $23.15 : 1$ (Tiết kiệm $95.68\%$ đĩa!)}
  \item Xử lý vector hóa Batch Mode (mảng $\sim 900$ dòng/lô trong CPU cache).
\end{itemize}
\end{exampleblock}
\end{column}

\begin{column}{0.48\textwidth}
\begin{alertblock}{OPEN Delta Rowgroup \& Nuance POINT}
\scriptsize
\begin{itemize}
  \item \textbf{Ghi chèn/sửa:} Chuyển vào OPEN Delta Rowgroup $\to$ throughput ghi giảm $-49.9\%$ (13,265 $\to$ 6,642 dòng/s). UPDATE ghi vào Delete Bitmap.
  \item \textbf{Bản chất POINT:} Columnstore scan segment nén nhanh hơn heap thô, nhưng B-tree seek thực sự chỉ mất $<1$ ms và 2--3 trang.
\end{itemize}
\end{alertblock}
\end{column}
\end{columns}
\end{frame}

% -----------------------------------------------------------------------------
% SLIDE 20: UC6 MULTIKEY VS NORMALIZED CHILD (SETUP)
% -----------------------------------------------------------------------------
\begin{frame}{UC6: Multikey vs. Normalized Child Index (10K Parents / 30K Tags)}
\textbf{Bài toán thực tế:} Mỗi đối tượng có thể gắn từ 3 đến 5 tags. Cần tìm các đối tượng chứa 1 tag hoặc đồng thời chứa 2 tags.

\begin{columns}[T]
\begin{column}{0.48\textwidth}
\begin{block}{MongoDB 7.0 (Mảng Nhúng Multikey)}
\footnotesize
\begin{itemize}
  \item Cấu trúc: Document chứa mảng \texttt{tags: ["type:TABLE", "owner:SYS"]}.
  \item Multikey Index trên \texttt{tags}: Tự động unnest mảng và tạo B-tree entry cho từng phần tử tag.
  \item Lưu trữ: 1 document duy nhất, không cần JOIN.
\end{itemize}
\end{block}
\end{column}

\begin{column}{0.48\textwidth}
\begin{block}{SQL Server 2025 (Bảng Con Chuẩn Hóa 1NF)}
\footnotesize
\begin{itemize}
  \item Cấu trúc: Bảng cha \texttt{Parent} (10K dòng) và bảng con \texttt{Parent\_Tags} (30,471 dòng quan hệ).
  \item Composite B-Tree Index trên bảng con: \texttt{(TAG, PARENT\_ID)}.
  \item Truy vấn 2 tag: Dùng \texttt{INTERSECT} hoặc Gom nhóm \texttt{HAVING COUNT = 2}.
\end{itemize}
\end{block}
\end{column}
\end{columns}

\begin{exampleblock}{Bốn Câu Truy Vấn Kiểm Thử Hành Vi}
\footnotesize
\texttt{ONE\_TABLE} (1,247 đối tượng) $\cdot$ \texttt{TWO\_SYS\_TABLE} (1,125 đối tượng) $\cdot$ \texttt{ONE\_TEMP\_Y} (51 đối tượng) $\cdot$ \texttt{TWO\_TABLE\_TEMP\_Y} (41 đối tượng).
\end{exampleblock}
\end{frame}

% -----------------------------------------------------------------------------
% SLIDE 21: UC6 RESULTS & CANDIDATE FETCHING
% -----------------------------------------------------------------------------
\begin{frame}{UC6: Kết Quả Thực Nghiệm, Candidate Fetching \& Đánh Đổi}
\begin{table}
\centering\footnotesize
\begin{tabularx}{\linewidth}{@{}l l r r r X@{}}
\toprule
\textbf{Hệ Quản Trị} & \textbf{Truy Vấn Tag} & \textbf{Baseline A} & \textbf{Index C} & \textbf{Speedup} & \textbf{Actual Plan Execution (A $\to$ C)} \\
\midrule
SQL Server & \texttt{ONE\_TABLE} (1,247) & 13.12 ms & \textbf{4.00 ms} & \textbf{3.28$\times$} & $150 \to \mathbf{8\text{ trang}}$ (PK Scan $\to$ Seek bảng con) \\
SQL Server & \texttt{TWO\_SYS\_TABLE} (1,125) & 16.66 ms & \textbf{11.69 ms} & \textbf{1.42$\times$} & $150 \to \mathbf{41\text{ trang}}$ (2 Seeks + Stream Aggregate) \\
SQL Server & \texttt{ONE\_TEMP\_Y} (51) & 10.51 ms & \textbf{0.00$^*$ ms} & N/A$^*$ & $150 \to \mathbf{2\text{ trang}}$ (p95: 11.2 $\to$ 1.0 ms) \\
SQL Server & \texttt{TWO\_TABLE\_TEMP\_Y} (41) & 10.85 ms & \textbf{2.01 ms} & \textbf{5.39$\times$} & $150 \to \mathbf{10\text{ trang}}$ (2 Seeks + Stream Aggregate) \\
\midrule
MongoDB & \texttt{ONE\_TABLE} (1,247) & 18.08 ms & \textbf{11.58 ms} & \textbf{1.56$\times$} & 10,000 docs $\to \mathbf{1,247\text{ keys \& docs examined}}$ \\
MongoDB & \texttt{TWO\_SYS\_TABLE} (1,125) & 21.73 ms & \textbf{13.07 ms} & \textbf{1.66$\times$} & Seek tag TABLE $\to$ FETCH 1,125 candidate docs \\
MongoDB & \texttt{ONE\_TEMP\_Y} (51) & 28.24 ms & \textbf{2.84 ms} & \textbf{9.94$\times$} & 10,000 docs $\to \mathbf{51\text{ keys \& docs examined}}$ \\
MongoDB & \texttt{TWO\_TABLE\_TEMP\_Y} (41) & 28.53 ms & \textbf{2.06 ms} & \textbf{13.85$\times$} & Seek tag TEMP\_Y (51 keys) $\to$ FETCH filter RAM \\
\bottomrule
\end{tabularx}
\end{table}

\vspace{-0.2cm}
\begin{columns}[T]
\begin{column}{0.48\textwidth}
\begin{block}{Chiến Lược Tối Ưu Của MongoDB}
\scriptsize
Trên truy vấn 2 tag, multikey entry chỉ lưu 1 tag. MongoDB chọn seek tag có tính chọn lọc cao nhất (\texttt{temporary:Y} chỉ có 51 keys), sau đó \textbf{FETCH 51 candidate docs} để kiểm tra tag thứ hai trong bộ nhớ RAM $\to$ \textbf{Speedup 13.85x}!
\end{block}
\end{column}

\begin{column}{0.48\textwidth}
\begin{alertblock}{Đánh Đổi Ghi Mảng Đa Trị (Table B)}
\scriptsize
\begin{itemize}
  \item \textbf{Nhân bản Index Entry:} Trung bình 3.05 tag/doc $\to$ tạo ra 30,471 entries index cho 10,000 document.
  \item \textbf{Giảm thông lượng ghi:} INSERT trên Mongo giảm \textbf{$-29.3\%$}; thao tác thêm tag vào mảng (\texttt{\$addToSet}) tăng \textbf{$+33.0\%$} độ trễ ghi.
\end{itemize}
\end{alertblock}
\end{column}
\end{columns}
\end{frame}

% -----------------------------------------------------------------------------
% SLIDE 22: SYNTHESIS - ANSWERING CORE QUESTIONS
% -----------------------------------------------------------------------------
\begin{frame}{Chương 5: Trả Lời Trực Tiếp Ba Câu Hỏi Cốt Lõi Của Đề Bài}
\begin{enumerate}
  \item \textbf{What are they? (Những cấu trúc/kỹ thuật nào được nghiên cứu?)}
  \begin{itemize}
    \item \emph{Kỹ thuật B-Tree:} Compound B-Tree (UC1), Filtered/Partial Index (UC2), Covering Index (UC3), Multikey Index (UC6).
    \item \emph{Cấu trúc vật lý khác:} Inverted Token Index cho Full-Text (UC4) và Compressed Columnar Segments cho Columnstore (UC5).
  \end{itemize}
  \item \textbf{How are they supported? (Các DBMS hiện thực và tối ưu ra sao?)}
  \begin{itemize}
    \item SQL Server hỗ trợ \texttt{INCLUDE} leaf-payload, Background Async Full-text Catalog, và Batch Mode vectorization trên Clustered Columnstore.
    \item MongoDB hỗ trợ Partial Filter Expressions, Multikey Array Unnesting, Synchronous Text Indexing và Covered Compound Projections (\texttt{\_id: 0}).
  \end{itemize}
  \item \textbf{How they support? (Chúng hỗ trợ cải thiện hiệu năng thế nào?)}
  \begin{itemize}
    \item Tăng tốc vượt trội từ $3\times$ đến $395\times$ khi truy vấn tận dụng đúng đường truy cập vật lý (seek đa cột, loại bỏ Sort, bỏ RID lookup, batch mode, nén đĩa $95\%$).
  \end{itemize}
\end{enumerate}
\end{frame}

% -----------------------------------------------------------------------------
% SLIDE 23: MASTER DECISION MATRIX
% -----------------------------------------------------------------------------
\begin{frame}{Bảng Quyết Định Kỹ Thuật Tổng Hợp (V3 PDF Bảng 5.1)}
\tiny
\begin{table}
\centering
\begin{tabularx}{\linewidth}{@{}l X X X@{}}
\toprule
\textbf{Chỉ Mục Nghiên Cứu} & \textbf{Lợi Ích Nổi Bật Do Được} & \textbf{Khi Nào NÊN DÙNG} & \textbf{Đánh Đổi / Khi Nào TRÁNH DÙNG} \\
\midrule
\textbf{UC1: Compound B-Tree} & U1a: 3.4x (SQL), 103x (Mongo); U1d: Triệt tiêu toán tử Sort. & Lọc nhiều cột luôn có leading key, hoặc ORDER BY khớp khóa. & Tránh khi thiếu leading key (U1c); Index lớn hơn 90\%, Mongo INSERT giảm 33.5\%. \\
\addlinespace
\textbf{UC2: Filtered / Partial} & Query 'Y' 1M: 15.4x (SQL), 395x (Mongo) so với A. & Bảng lệch nặng, thường xuyên lọc giá trị hiếm đã biết (0.24\% active). & Tránh khi query dòng ngoài tập ('N') hoặc tham số SQL (@p); Tiết kiệm $>98\%$ đĩa. \\
\addlinespace
\textbf{UC3: Covering Index} & Medium range: 1.54x (SQL 1M), 128x (Mongo vs A); 0 base lookups. & Range query tần suất cao trả về nhiều cột nơi lookup làm nghẽn I/O. & Tránh cho Point lookup (1 dòng); Dung lượng 1M phình to gấp 6x đến 10x! \\
\addlinespace
\textbf{UC4: Full-Text / Text} & SQL: 4s $\to$ 15--20 ms; Mongo: 3.3x--107.8x trên cùng tập ID. & Tìm kiếm từ và cụm từ ngôn ngữ tự nhiên bên trong văn bản dài. & Tránh khi cần substring chính xác; SQL trễ 5.8s hiển thị, Mongo INSERT giảm 95.3\%. \\
\addlinespace
\textbf{UC5: Clustered Columnstore} & AGG: 2.90x nhanh hơn (Batch Mode); Nén 23.15 : 1 (219 $\to$ 9.46 MB). & Kho dữ liệu, bảng báo cáo phân tích quét tổng hợp ít cột trên 1M dòng. & Tránh cho OLTP ghi nhỏ liên tục; Build mất 22s, INSERT giảm 50\% qua Delta Store. \\
\addlinespace
\textbf{UC6: Multikey / Child Index} & Tag chọn lọc: 3.3x--5.4x (SQL) và 9.9x--13.9x (Mongo) so với A. & Lọc theo mảng tags hoặc dòng bảng con theo tên tag chọn lọc. & Tránh cho mảng quá dài ghi nhiều; Lưu 3.05 entries/doc, Mongo INSERT giảm 29.3\%. \\
\bottomrule
\end{tabularx}
\end{table}
\end{frame}

% -----------------------------------------------------------------------------
% SLIDE 24: METHODOLOGICAL PITFALLS
% -----------------------------------------------------------------------------
\begin{frame}{Bốn Cạm Bẫy Phương Pháp Luận Cần Tránh Khi Đánh Giá Index}
\begin{enumerate}
  \item \textbf{Có index chưa chắc đã nhanh hơn quét bảng (Index Scan vs Table Scan):} Trong UC1 Mongo, đánh index đơn cột trên \texttt{OWNER='SYS'} làm chậm hơn scan gần gấp đôi (462.8 ms vs 241.0 ms) do duyệt 71,802 key và FETCH 71,802 doc ngẫu nhiên!
  \item \textbf{Thời gian chạy thấp chưa chắc là Speedup cùng kết quả (Semantic Parity):} Trong UC4 SQL, \texttt{LIKE} và \texttt{CONTAINS} trả về khác tập ID ở cả 6 query do thuật toán tách từ và stemming. Ta không được gán ghép tỉ lệ tăng tốc khi câu hỏi ngữ nghĩa khác nhau.
  \item \textbf{Quét Columnstore nhanh hơn không có nghĩa là thay thế được B-Tree Seek:} Trong UC5, quét segment nén nhanh hơn quét trang heap ($3.63$ ms vs $88.8$ ms), nhưng đây vẫn là scan toàn bộ bảng chứ không thể đạt tốc độ dưới 1 ms và đọc 2 trang như B-tree seek.
  \item \textbf{Ngưỡng độ phân giải timer dưới mili-giây (0.00 ms Timer Floor):} Trên SQL Server ở các query selective (UC3 Point, UC6 ONE\_TEMP\_Y), timer batch ghi nhận 0.00 ms. Không được chia cho 0 để ra 'vô cực lần', mà phải đánh giá qua số trang đọc vật lý.
\end{enumerate}
\end{frame}

% -----------------------------------------------------------------------------
% SLIDE 25: STREAMLIT COMPANION DEMO ROADMAP
% -----------------------------------------------------------------------------
\begin{frame}{Kế Hoạch Live Demo Web Streamlit}
\begin{columns}[T]
\begin{column}{0.48\textwidth}
\begin{block}{Cấu Trúc Ứng Dụng Web Ba Phần}
\begin{enumerate}
  \item \textbf{Directory of Workload Queries:} Hiển thị tường minh từng query riêng lẻ cho mỗi UC (U1a--U1d, Y/N/@p, 3 Selective Tiers, 6 FTS terms, AGG/POINT, 4 Tags).
  \item \textbf{Interactive Benchmark Visualizer:} Biểu đồ Plotly nhóm so sánh A vs B vs C; lọc động theo DBMS, quy mô dữ liệu và chỉ số p50/p95/IQR.
  \item \textbf{Engine Architecture \& Execution Plans:} Trực quan hóa Actual Plans, toán tử vật lý, bảng đánh đổi chi phí ghi (Table B) và ma trận quyết định.
\end{enumerate}
\end{block}
\end{column}

\begin{column}{0.48\textwidth}
\begin{exampleblock}{Thông Số Trực Quan Trọng Tâm Sẽ Demo}
\begin{itemize}
  \item \textbf{UC1:} Minh họa hiện tượng sụp đổ của Mongo B trên đồ thị và biểu diễn loại bỏ toán tử Sort.
  \item \textbf{UC2:} Trực quan hóa tỉ lệ nén đĩa $99.71\%$ bằng biểu đồ cột dung lượng (0.13 MB vs 43.74 MB).
  \item \textbf{UC3:} Điểm bùng nổ Tipping Point khi vượt qua 1,600 dòng kết quả.
  \item \textbf{UC4:} Bảng đối soát tập ID văn bản và thước đo trễ ghi nhận $5.8$s của SQL Server.
  \item \textbf{UC5:} Tỉ lệ nén $23.15:1$ và cơ chế Delta Store.
\end{itemize}
\end{exampleblock}
\end{column}
\end{columns}
\end{frame}

% -----------------------------------------------------------------------------
% SLIDE 26: CONCLUSION & Q&A
% -----------------------------------------------------------------------------
\begin{frame}
\centering
\vspace{0.8cm}
{\Huge\textbf{\color{hcmutblue}XIN CHÂN THÀNH CẢM ƠN}}\\[0.3cm]
{\Large\textbf{THẦY VÀ CÁC BẠN ĐÃ LẮNG NGHE!}}\\[0.8cm]

\begin{alertblock}{Thông Điệp Kỹ Sư Cốt Lõi Từ Đồ Án}
\centering
\textbf{Không có ``chỉ mục vạn năng'' (Silver Bullet Index).} Mọi chỉ mục đều là sự đánh đổi cân não giữa \textbf{chi phí đĩa, thời gian build và thông lượng ghi} để đổi lấy \textbf{tốc độ đọc}. Hãy luôn đo đạc bằng Workload thực tế và giám sát chặt chẽ qua Actual Execution Plan!
\end{alertblock}

\vspace{0.5cm}
{\large\textbf{\color{accentblue}Nhóm 02 sẵn sàng trả lời các câu hỏi (Q\&A)}}\\[0.2cm]
\footnotesize CO3021 - Hệ Quản Trị Cơ Sở Dữ Liệu $\cdot$ Học kỳ 261 $\cdot$ HCMUT 2026
\end{frame}

\end{document}
